% ----------------------------------------------------------
\begin{frame}[fragile]{Subqueries -- Introduction}
  A \textbf{subquery} (inner query) is a \texttt{SELECT} statement nested inside
  another SQL statement. It is evaluated first and its result is used by the outer query.

  \begin{block}{Three Positions for Subqueries}
    \begin{tabular}{@{}ll@{}}
      \toprule
      \textbf{Position} & \textbf{What it produces} \\
      \midrule
      \texttt{FROM} clause   & A temporary table (derived table / inline view) \\
      \texttt{WHERE} clause  & A value or set of values to compare against \\
      \texttt{SELECT} clause & A scalar value (one value per outer row) \\
      \bottomrule
    \end{tabular}
  \end{block}

  \begin{block}{Example: Subquery in FROM}
\begin{lstlisting}[style=sqlstyle]
-- Find artist name for albums released before 1980
SELECT ar.Name, sub.Title
FROM (
  SELECT Title, ArtistID FROM Album WHERE ReleaseYear < 1980
) AS sub
JOIN Artist AS ar ON sub.ArtistID = ar.ArtistID;
-- Returns: The Beatles / Abbey Road, Kraftwerk / Autobahn
\end{lstlisting}
  \end{block}
\end{frame}

% ----------------------------------------------------------
\begin{frame}[fragile]{Subqueries -- IN and EXISTS}
  \begin{block}{Subquery in WHERE with IN}
\begin{lstlisting}[style=sqlstyle]
-- Artists who have at least one Electronic album
SELECT Name FROM Artist
WHERE ArtistID IN (
  SELECT ArtistID FROM Album WHERE Genre = 'Electronic'
);
-- Returns: Kraftwerk, Daft Punk
\end{lstlisting}
  \end{block}

  \begin{block}{Subquery with EXISTS (often faster)}
\begin{lstlisting}[style=sqlstyle]
-- Same result, but EXISTS stops as soon as one match is found
SELECT Name FROM Artist AS ar
WHERE EXISTS (
  SELECT 1 FROM Album AS al
  WHERE al.ArtistID = ar.ArtistID
    AND al.Genre = 'Electronic'
);
\end{lstlisting}
  \end{block}

  \begin{itemize}
    \item \texttt{EXISTS} returns \texttt{TRUE} as soon as one matching row is found
    \item \texttt{NOT EXISTS} finds artists with \emph{no} Electronic albums
  \end{itemize}
\end{frame}

% ----------------------------------------------------------
\begin{frame}[fragile]{Correlated Subqueries -- and Why to Avoid Them}
  \begin{block}{Correlated subquery: re-executed once per outer row}
\begin{lstlisting}[style=sqlstyle]
-- For each track, show the average duration of its album
SELECT t.Title, t.Duration,
  (SELECT AVG(t2.Duration) FROM Track t2
   WHERE t2.AlbumID = t.AlbumID) AS album_avg
FROM Track AS t;
\end{lstlisting}
  \end{block}

  \begin{block}{Better: JOIN + GROUP BY (executes only once)}
\begin{lstlisting}[style=sqlstyle]
SELECT t.Title, t.Duration, avg_data.avg_dur
FROM Track AS t
JOIN (
  SELECT AlbumID, AVG(Duration) AS avg_dur
  FROM Track GROUP BY AlbumID
) AS avg_data ON t.AlbumID = avg_data.AlbumID;
\end{lstlisting}
  \end{block}

  \begin{alertblock}{Performance}
    A correlated subquery re-executes \textbf{once per outer row}.
    With 50,000 tracks that is 50,000 subquery executions. Always prefer JOIN.
  \end{alertblock}
\end{frame}

% ----------------------------------------------------------
\begin{frame}[fragile]{EXISTS, ANY, ALL}
  \begin{block}{Set-Based Comparisons}
\begin{lstlisting}[style=sqlstyle]
-- ANY: true if ANY value in subquery satisfies the condition
-- Tracks longer than AT LEAST ONE track on Album 1
SELECT Title, Duration FROM Track
WHERE Duration > ANY (
  SELECT Duration FROM Track WHERE AlbumID = 1
);
-- Equivalent to: Duration > MIN(durations on Album 1) = 182

-- ALL: true if the condition holds for ALL subquery values
-- Tracks longer than EVERY track on Album 1
SELECT Title, Duration FROM Track
WHERE Duration > ALL (
  SELECT Duration FROM Track WHERE AlbumID = 1
);
-- Equivalent to: Duration > MAX(durations on Album 1) = 259
\end{lstlisting}
  \end{block}

  \begin{itemize}
    \item \texttt{> ANY} $\equiv$ \texttt{> MIN(\ldots)}
    \item \texttt{> ALL} $\equiv$ \texttt{> MAX(\ldots)}
    \item \texttt{NOT EXISTS} is the preferred way to express relational division
  \end{itemize}
\end{frame}
